% Pursuing Stacks (cotes 134-2 to 134-8): archive cote 134-4, site page 46 (PDF page 47), typescript page 134.
% typescript. Draft compared: ps3.tex:2332-2398.
% AI-transcribed by Claude Opus 5.5 (claude-opus-5-5, adaptive thinking, effort high), Message Batches API, batch ps2 (msgbatch_01NVwr1Yuo8qJf5rykeGBXPK), 2026-11-12; not yet checked by eye.
\scanpage{134-4/46}
\mspage{134}
also the intervals such that $I \hand{\to} e$ is \xout{\unsure{xxwxxk}} in W (and hence universally
so in terms of the assumption made of compatibility of u \over{on}{with} products).
The trouble with \xout{this} working with this h,  rather than with $h_W$ as above,
is twofold though:  a)  The assumptions of strong saturation on W and com-%
patibility with products \xout{\unsure{inxnot}} are not so readily verified in the cases
of interest to us,  and the second moreover is not always satisfied,  e.g.
there are test categories which are not strict i.e. elementary modelizers
which are not strict;  b)  the condition $u(\hand{\delta}_0) = u(\hand{\delta}_1)$  is not readily
verified in terms of W directly,  whereas the condition of W-asphericity
is  -  still more so if W is compatible with products and hence $I \hand{\to} e$ is
W-asperic just means it is in W.

55.     Let's now look at ``\uline{morphisms}'' between categories endowed with
homotopy structures,  \xout{(M,h)xxandx(M',h')xxsayx}  $(M_1,h_1)$ and $(M_2,h_2)$ say.
The natural definition here is to take as morphisms between these homotopy
structures the functors   $u\colon M_1 \hand{\to} M_2$  which give rise to a commutative
square of functors
\diagram{\begin{tikzcd}
M_1 \ar[r,"u"] \ar[d] & M_2 \ar[d] \\
\overline{M}_1 \ar[r,"\overline{u}"] & \overline{M}_2
\end{tikzcd}\quad ,}
where the vertical arrows are the canonical functors into the respective
homotopy-types categories,  and $\overline{u}$  a suitable functor,  necessarily unique.
The existence of $\overline{u}$  can be expressed at will in terms of the $\mathrm{Hom}_1$ or $\mathrm{Hom}_2$
structures,  namely as
\[ \text{(i)}\quad f \underset{h_1}{\hand{\sim}} g \quad\text{implies}\quad u(f) \underset{h_2}{\hand{\sim}} u(g) \quad , \]
\over{and}{or} \hand{as}
\[ \text{(ii)}\quad f \hand{\in} W_{h_1} \quad\text{implies}\quad u(f) \hand{\in} W_{h_2} \quad . \]
These conditions,  when $M_1$ and $M_2$ have final objects and these are respec\xout{t}-%
ted by u,  imply that u transforms \inserted{weak} homotopy intervals into weak homotopy
intervals,  and homotopisms into homotopisms.  Conversely,  \uline{if} u \uline{commutes
with finite products},  and if $h_1$ \over{\illegible{1}}{c}an be defined by a weak homotopy interv\hand{al}
structure (respectively,  by a contractibility structure),  than for u to
be a morphism of homotopy structures,  it is (necessary and) sufficient
that u carry weak homotopy intervals (resp. contractible objects) into
same.

In case $h_1$ is described in terms of a generating set $\Sigma^{0}_{\hand{h}_1}$  of weak
homotopy intervals,  and \inserted{if} u \uline{commutes with finite products},  the \xout{\illegible{1}}
most economic way often to express that u is a morphism of homotopy struc-%
tures,  is by the condition that  for any $\underline{I}$ in $\Sigma^{0}_{h_1}$ ,  $u(\underline{h})$ be a weak ho-%

\margin{46}
